\documentclass[12pt]{article}

% ── Layout ──────────────────────────────────────────────────────────────────
\usepackage[margin=1in]{geometry}
\usepackage{setspace}
\doublespacing

% ── Math ────────────────────────────────────────────────────────────────────
\usepackage{amsmath,amssymb,amsthm}

% ── Tables ──────────────────────────────────────────────────────────────────
\usepackage{booktabs}
\usepackage{multirow}
\usepackage{array}

% ── Fonts and typography ────────────────────────────────────────────────────
\usepackage{mathptmx}
\usepackage{microtype}

% ── References ──────────────────────────────────────────────────────────────
\usepackage[authoryear,round]{natbib}
\usepackage[colorlinks=true,
            linkcolor=black,
            citecolor=black,
            urlcolor=black]{hyperref}

% ── Misc ────────────────────────────────────────────────────────────────────
\usepackage{caption}
\usepackage{subcaption}

% ─────────────────────────────────────────────────────────────────────────────

\begin{document}

% ════════════════════════════════════════════════════════════════════════════
%  TITLE PAGE
% ════════════════════════════════════════════════════════════════════════════

\begin{titlepage}
\centering
\vspace*{1.5cm}

{\LARGE\bfseries Gradient-Boosted Trees for Realized Volatility Forecasting:\\[8pt]
Fractional Differentiation, Domain Constraints,\\[8pt]
and the Limits of Stacking\par}

\vspace{1.5cm}

{\large [Author Name]\par}
\vspace{0.4cm}
{\normalsize MS Quantitative Finance and Risk Management\\
Fox School of Business, Temple University\par}

\vspace{1.2cm}

{\normalsize April 2026\par}

\vspace{2cm}

\begin{abstract}
\noindent
This paper evaluates a suite of gradient-boosted tree models against the
GARCH(1,1)-$t$ benchmark for next-day five-day realized volatility forecasting
on the SPDR S\&P 500 ETF (SPY) using ten years of daily data spanning 2014
through 2024.  Three methodological contributions are examined alongside the
standard ML pipeline: fractional differentiation of log-prices at the minimum
integration order consistent with stationarity ($d^{*} = 0.6$ for SPY, ADF
$p = 0.0007$), which preserves long-memory properties destroyed by integer
differencing; a monotone-constrained XGBoost model that encodes the leverage
effect and volatility clustering as hard architectural constraints rather than
soft regularizers; and a leakage-aware stacking ensemble using purged,
embargoed walk-forward cross-validation aligned with the five-day target
horizon.  On a holdout test of 485 trading days from January 2023 through
December 2024, XGBoost reduces RMSE by 16.4\% relative to GARCH (0.002334
vs.\ 0.002793).  The quantile median estimator (Q50) achieves the best
performance on all four reported metrics: RMSE = 0.002156, QLIKE $= -8.8054$,
MAPE = 25.62\%.  Stacking ensembles fail on every metric, underperforming
each of their three component models by substantial margins.  The
monotone-constrained model incurs a QLIKE cost of 0.032 relative to
unconstrained XGBoost in exchange for full financial interpretability.  The
results suggest that robust estimation of the conditional median, combined
with a rich but carefully debiased feature set, outperforms both classical
econometrics and naive ensemble aggregation for this forecasting task.
\end{abstract}

\vspace{1cm}

\noindent\textbf{JEL Classification:} C22, C53, C63, G17

\vspace{0.3cm}

\noindent\textbf{Keywords:} realized volatility, GARCH, gradient boosting,
fractional differentiation, QLIKE loss, monotone constraints, quantile
regression, stacking ensemble

\end{titlepage}

\newpage

% ════════════════════════════════════════════════════════════════════════════
%  1. INTRODUCTION
% ════════════════════════════════════════════════════════════════════════════

\section{Introduction}

Accurate volatility forecasts are inputs to portfolio risk management, option
pricing, and regulatory capital allocation.  The question of whether machine
learning can improve on the GARCH(1,1) benchmark has motivated a large
literature, yet the answer depends critically on three often-ignored factors:
the construction of the baseline, the leakage-awareness of the cross-validation
scheme, and the choice of evaluation metric.  This paper addresses all three
while applying gradient-boosted trees to daily realized volatility forecasting
on US equity data.

The GARCH(1,1) model of \citet{Bollerslev1986} remains a formidable benchmark.
\citet{HansenLunde2005} compared 330 volatility models across a large currency
dataset and found no model that consistently beats GARCH(1,1) in pairwise
comparisons.  Part of this durability comes from parsimony: GARCH captures
volatility clustering with three parameters and generalizes reliably
out-of-sample.  Any ML model claiming to improve on GARCH must earn that
improvement honestly, on a holdout set, under a metric that is robust to the
noisiness of the realized volatility proxy.

This paper makes four specific contributions beyond applying gradient boosting
to volatility data.

First, the GARCH(1,1)-$t$ baseline is constructed to use only the
one-step-ahead conditional variance forecast $\hat{\sigma}_{t+1|t}^{2}$,
which depends solely on information available at the close of day $t$.  A
common implementation error lags the GARCH output series incorrectly, allowing
same-day shock information to enter the baseline and flatter its
out-of-sample performance.

Second, the log-price series is fractionally differenced at the minimum
integration order achieving stationarity on the training fold ($d^{*} = 0.6$
for SPY, found by ADF grid search).  Fractional differencing preserves the
hyperbolic decay of autocorrelations that characterizes financial price series
\citep{GrangerJoyeux1980}, whereas first-differencing over-removes this
dependence.  The resulting feature retains memory of price levels from over a
year ago while satisfying the stationarity requirement for tree-based models.

Third, an XGBoost model encodes the leverage effect \citep{Black1976} and
volatility clustering as hard monotone constraints embedded in the tree-split
algorithm.  This converts domain knowledge from a regularization prior into an
architectural guarantee: the model cannot produce a prediction that contradicts
the sign of these well-established economic relationships, regardless of the
input configuration.

Fourth, the stacking ensemble uses manually generated out-of-fold predictions
with a five-day embargo at each fold boundary.  The embargo prevents the
five-day realized volatility target at the boundary from contaminating the
base model's training data through shared return observations.

Results on 485 out-of-sample trading days show that XGBoost improves RMSE by
16.4\% over GARCH, the quantile median estimator achieves the best performance
on all reported metrics, and stacking ensembles underperform every individual
base model.  Section~\ref{sec:lit} reviews the relevant literature.
Section~\ref{sec:data} describes the data.  Section~\ref{sec:method} presents
the methodology.  Section~\ref{sec:results} reports results.
Section~\ref{sec:disc} discusses findings and limitations.
Section~\ref{sec:conc} concludes.

% ════════════════════════════════════════════════════════════════════════════
%  2. LITERATURE REVIEW
% ════════════════════════════════════════════════════════════════════════════

\section{Literature Review}
\label{sec:lit}

\subsection*{Econometric Volatility Models}

\citet{Engle1982} introduced the ARCH model, formalizing the observation that
return variances cluster: large shocks are followed by large shocks regardless
of sign.  \citet{Bollerslev1986} extended this to the GARCH(p,q)
specification, which allows conditional variance to depend on lagged squared
residuals and lagged conditional variances.  The GARCH(1,1) parameterization
captures most empirically observed volatility dynamics parsimoniously.
\citet{Bollerslev1987} replaced the normal innovation distribution with the
Student-$t$, providing better in-sample fit for the heavy tails characteristic
of equity returns.  These three papers form the econometric baseline against
which ML is benchmarked here.

\subsection*{Realized Volatility and HAR-RV}

The availability of high-frequency data motivated measuring realized volatility
directly as the sum of squared intraday returns.  \citet{AndersenEtAl2003}
established the theoretical and empirical basis for realized variance as a
consistent estimator of integrated variance and demonstrated that realized
measures produce substantially more accurate forecasts than GARCH-family models
when used as forecast inputs.  \citet{Corsi2009} proposed the Heterogeneous
Autoregressive model for Realized Volatility (HAR-RV), which regresses
next-period realized volatility on daily, weekly, and monthly averages of past
realized volatility, capturing the distinct trading horizons of different market
participants.  Despite its simplicity, HAR-RV consistently outperforms
GARCH(1,1) when realized measures are available \citep{HansenLunde2005}.  The
present study implements HAR-RV in the codebase but does not compare it
numerically in the main results, a limitation addressed in
Section~\ref{sec:disc}.

\subsection*{High-Efficiency Volatility Estimators from OHLC Data}

When intraday data are unavailable, efficient volatility estimates can be
extracted from daily open-high-low-close bars.  \citet{Parkinson1980} showed
that the price range contains more information about the day's volatility than
the close-to-close return alone.  \citet{GarmanKlass1980} derived an optimal
estimator using the full OHLC bar, achieving approximately seven times lower
variance than the close-to-close estimator under geometric Brownian motion.
Both are used as features in this paper.

\subsection*{The Leverage Effect}

\citet{Black1976} documented that negative equity returns are followed by
disproportionately large increases in volatility.  The standard interpretation
is that a falling stock price raises the firm's debt-to-equity ratio and
therefore its equity risk.  This asymmetric response is encoded directly into
the Entropy XGBoost model as a monotone constraint: the model's prediction
cannot increase when lagged log returns increase, holding other features fixed.

\subsection*{Long Memory and Fractional Integration}

\citet{GrangerJoyeux1980} introduced the concept of fractional integration.
For a process integrated of order $d \in (0,1)$, the autocorrelation function
decays hyperbolically as $\rho(k) \sim k^{2d-1}$, slower than any
exponentially declining ARMA process.  First-differencing such a process
$(d = 1)$ over-differences it, destroying the long-range dependence.
\citet{LopezDePrado2018} applied fractional differencing to financial prices
as a feature engineering technique: choosing the minimum $d$ that achieves
stationarity preserves the maximum amount of economically relevant memory
while satisfying the stationarity requirement that machine learning models
implicitly assume.

\subsection*{Evaluation Metrics for Volatility Forecasts}

A core difficulty in comparing volatility forecasts is that the target, true
latent conditional variance, is unobservable.  \citet{Patton2011} showed that
not all loss functions produce rankings that are invariant to the quality of the
realized-volatility proxy used in place of latent variance.  Theorem~1 of
\citet{Patton2011} establishes that the QLIKE loss produces consistent model
rankings under a broad class of proxies, while mean squared error (MSE)
rankings can be reversed by proxy noise.  This finding motivates QLIKE as the
primary evaluation metric in this paper.  The \citet{DieboldMariano1995} test
for equal predictive accuracy is implemented in the accompanying codebase using
Newey-West long-run variance to account for the serial correlation induced by
the overlapping five-day target.

\subsection*{Gradient Boosting}

\citet{Friedman2001} formalized gradient boosting as functional gradient
descent in a space of weak learners, providing the theoretical foundation for
the tree-ensemble methods used here.  \citet{ChenGuestrin2016} introduced
XGBoost, which added approximate split-finding, column subsampling, and
$L_2$ regularization, enabling reliable large-scale applications.
\citet{KeMengEtAl2017} developed LightGBM with leaf-wise tree growth,
producing models with more leaves and different effective depths than
depth-bounded XGBoost at comparable computational cost.  Both implementations
support monotone constraints on individual features, where a proposed split is
rejected if it violates the specified partial-monotonicity requirement.

\subsection*{Ensemble Methods and Stacking}

\citet{RapachZhou2013} surveyed equity return forecasting and found that
forecast combination frequently outperforms individual models, but that this
advantage depends on the diversity of the constituent forecasts.  Stacking
trains a second-level model on out-of-fold predictions of base models and
requires careful handling of temporal structure to avoid in-sample contamination.
This paper finds that stacking underperforms every base model in the volatility
setting, consistent with the intuition that highly correlated base forecasts
provide little material for a meta-learner to exploit.

% ════════════════════════════════════════════════════════════════════════════
%  3. DATA
% ════════════════════════════════════════════════════════════════════════════

\section{Data}
\label{sec:data}

\subsection*{Primary Asset}

The primary asset is the SPDR S\&P~500 ETF (SPY), which tracks the
S\&P~500 index and is among the most liquid exchange-traded products in the
world.  Daily open, high, low, adjusted close, and volume data are obtained
from Yahoo Finance via the \texttt{yfinance} library for the period January
2014 through December 2024, yielding approximately 2,510 raw trading days.
After constructing features with lookback windows and dropping observations
with missing values, 2,421 complete observations remain.

\subsection*{Cross-Asset Signals}

Three cross-asset series are included to provide macro and credit market
information.  The CBOE Volatility Index (VIX) measures the 30-day
risk-neutral variance of the S\&P~500 implied by options prices and is widely
used as a proxy for investor uncertainty.  The CBOE 3-Month Volatility Index
(VIX3M) covers a 93-day horizon; its ratio to VIX captures the slope of the
implied volatility term structure and provides information about the expected
path of volatility beyond what the VIX level alone conveys.  Credit market
stress is proxied by the ratio of the iShares iBoxx USD High Yield Corporate
Bond ETF (HYG) to the iShares 7--10 Year Treasury Bond ETF (IEF): a declining
ratio indicates widening credit spreads and rising systemic financial stress.
All cross-asset series are sourced from Yahoo Finance.

\subsection*{Sample Split}

The sample is divided chronologically in an 80/20 ratio into a training set
and a test set.  The training set spans May 2015 through January 2023 and
contains 1,936 observations.  The test set spans January 2023 through December
2024 and contains 485 observations.  No shuffling is applied at any stage.
All preprocessing objects, including the feature scaler, the fractional
differencing order, and the volatility regime model, are estimated exclusively
on the training set and applied without refitting to the test set.

% ════════════════════════════════════════════════════════════════════════════
%  4. METHODOLOGY
% ════════════════════════════════════════════════════════════════════════════

\section{Methodology}
\label{sec:method}

\subsection{Target Variable}

The forecast target is the five-day realized volatility:
\begin{equation}
  \text{RV5}_{t} = \left[\frac{1}{4}\sum_{j=0}^{4}
    \bigl(r_{t-j} - \bar{r}_{t}\bigr)^{2}\right]^{1/2},
  \label{eq:rv5}
\end{equation}
where $r_{t} = \ln(C_{t}/C_{t-1})$ is the daily log return and
$\bar{r}_{t}$ is the mean of the five returns in the window.  The prediction
objective at each day $t$ is $\text{RV5}_{t+1}$, the realized volatility
realized over the five trading days beginning the following day.

Because RV5 is strictly positive, all ML models are trained on
$\ln(\text{RV5}_{t+1})$ and their outputs are exponentiated before evaluation:
\begin{equation}
  \hat{h}_{t+1} = \exp\!\bigl(\hat{f}(\mathbf{x}_{t})\bigr).
  \label{eq:logvol}
\end{equation}
The log transformation guarantees positive forecasts without constrained
optimization.

\subsection{Volatility Estimators}

Four daily volatility estimates computed from OHLC data serve as features.
Let $O_{t}$, $H_{t}$, $L_{t}$, $C_{t}$ denote the open, high, low, and
close prices on day $t$.

\smallskip\noindent\textit{Log return:}
$r_{t} = \ln(C_{t}/C_{t-1})$.

\smallskip\noindent\textit{Garman-Klass estimator} \citep{GarmanKlass1980}:
\begin{equation}
  \hat{\sigma}_{t}^{\,GK} =
    \sqrt{\frac{1}{2}\ln^{2}\!\left(\frac{H_{t}}{L_{t}}\right)
          - (2\ln 2 - 1)\ln^{2}\!\left(\frac{C_{t}}{O_{t}}\right)}.
  \label{eq:gk}
\end{equation}
This estimator uses the full OHLC bar and achieves approximately seven times
lower variance than the close-to-close estimator under geometric Brownian
motion.

\smallskip\noindent\textit{Parkinson estimator} \citep{Parkinson1980}:
\begin{equation}
  \hat{\sigma}_{t}^{\,P} =
    \sqrt{\frac{\ln^{2}(H_{t}/L_{t})}{4\ln 2}}.
  \label{eq:park}
\end{equation}

\smallskip\noindent\textit{Overnight gap:}
\begin{equation}
  \text{Gap}_{t} = \left|\ln\!\left(\frac{O_{t}}{C_{t-1}}\right)\right|.
  \label{eq:gap}
\end{equation}
The overnight gap captures discrete price jumps and after-hours news that do
not appear in intraday estimators.

\subsection{GARCH(1,1)-\texorpdfstring{$t$}{t} Baseline}
\label{sec:garch}

The GARCH(1,1) model with Student-$t$ innovations is:
\begin{align}
  r_{t} &= \mu + \varepsilon_{t}, \quad
  \varepsilon_{t} = \sigma_{t}\,z_{t}, \quad
  z_{t} \overset{iid}{\sim} t_{\nu}, \label{eq:garch-mean}\\
  \sigma_{t}^{2} &= \omega + \alpha\,\varepsilon_{t-1}^{2}
                   + \beta\,\sigma_{t-1}^{2}, \label{eq:garch-var}
\end{align}
with $\omega > 0$, $\alpha, \beta \geq 0$, and $\alpha + \beta < 1$.
Parameters are estimated by quasi-maximum likelihood on an expanding window,
refit every 21 trading days (approximately one calendar month).  The
one-step-ahead conditional volatility forecast used in the comparison is:
\begin{equation}
  \hat{\sigma}_{t+1|t} = \sqrt{\hat{\omega}
    + \hat{\alpha}\,\varepsilon_{t}^{2}
    + \hat{\beta}\,\hat{\sigma}_{t}^{2}},
  \label{eq:garch-forecast}
\end{equation}
which depends only on information available at the close of day $t$.
Concretely, the output series is lagged by one period before entering the
comparison, ensuring that the same-day residual $\varepsilon_{t}$ does not
contribute to the baseline forecast for $\text{RV5}_{t+1}$.  This is a
common but silent error in applied ML-versus-GARCH comparisons; correcting it
makes the baseline stronger.

\subsection{Feature Engineering}

\subsubsection*{Fractional Differentiation}

For $d \in (0,1)$, the fractionally differenced series is defined by the
binomial expansion:
\begin{equation}
  (1-L)^{d}\,x_{t} = \sum_{k=0}^{\infty} w_{k}\,x_{t-k},
  \label{eq:fracdiff}
\end{equation}
where the weights satisfy the recurrence $w_{0} = 1$ and
$w_{k} = -w_{k-1}(d - k + 1)/k$ for $k \geq 1$, truncated once
$|w_{k}| < 10^{-5}$.  For integer $d = 1$, the weights are $1, -1, 0, 0,
\ldots$ (finite support), destroying all memory beyond one lag.  For
$d \in (0,1)$, the weights decay as $|w_{k}| \sim k^{-(d+1)}$ (hyperbolic),
retaining influence from observations many periods in the past while reducing
the order of integration below one.

The minimum $d$ achieving stationarity (ADF $p \leq 0.05$) is found by grid
search over $d \in \{0.1, 0.2, \ldots, 1.0\}$ applied exclusively to
log-prices in the training fold.  For SPY, the result is $d^{*} = 0.6$ with
ADF $p = 0.0007$.  At this order, the weight at lag $k = 100$ is
$100^{-1.6} \approx 0.025$, so price information from more than four calendar
months ago still enters the feature.  Integer differencing would zero out all
such lags.

\subsubsection*{Volatility Regime Dummies}

A Gaussian Mixture Model (GMM) is estimated on the joint distribution of
$(\text{RV5}_{t}, \text{VIX\textrm{-}daily}_{t})$ using training data only,
where $\text{VIX\textrm{-}daily}_{t} = \text{VIX}_{t}/(100\sqrt{252})$
converts the annualized implied volatility index to a daily decimal scale.
The number of components is selected by the Bayesian Information Criterion
over $\{2, 3, 4, 5, 6, 7\}$; the BIC selects two components.  Regime labels
are sorted by ascending mean RV5, so regime 0 corresponds to low-volatility
conditions.  One-hot dummies for each regime are included as features and
provide a nonlinear partition of the volatility state space that linear
features cannot capture.

\subsubsection*{Additional Features}

Table~\ref{tab:features} summarizes the complete 18-variable feature set.
Cross-asset features capture macro and credit market information.  The
variance risk premium (VRP) is:
\begin{equation}
  \text{VRP}_{t} = \hat{\sigma}_{t+1|t}^{\,GARCH}
                  - \frac{\text{VIX}_{t}}{100\sqrt{252}},
  \label{eq:vrp}
\end{equation}
the gap between the GARCH conditional volatility forecast and the
options-implied daily volatility.  A positive VRP indicates that the options
market is pricing in more uncertainty than the GARCH model predicts based on
historical returns, consistent with a premium for bearing variance risk.

\begin{table}[htbp]
\centering
\caption{Feature Set ($p = 18$ base features; $p \times L = 360$ after
         lookback flattening with $L = 20$)}
\label{tab:features}
\begin{tabular}{llp{6.8cm}}
\toprule
Category & Feature & Construction \\
\midrule
\multirow{4}{*}{Vol.\ estimators}
  & Log return      & $r_{t} = \ln(C_{t}/C_{t-1})$ \\
  & Garman-Klass    & Equation~\eqref{eq:gk} \\
  & Parkinson       & Equation~\eqref{eq:park} \\
  & Overnight gap   & Equation~\eqref{eq:gap} \\
\midrule
\multirow{6}{*}{Technical}
  & RSI             & 14-period Relative Strength Index \\
  & BB Width        & $4 \times \text{sd}_{20}(r_{t})$ \\
  & Return Mom.     & $\text{mean}_{5}(r_{t})$ \\
  & Vol-of-Vol      & $\text{sd}_{20}(\hat{\sigma}_{t}^{GK})$ \\
  & Skewness        & $\text{skew}_{60}(r_{t})$ \\
  & GARCH Vol       & $\hat{\sigma}_{t+1|t}^{GARCH}$ \\
\midrule
\multirow{4}{*}{Cross-asset}
  & VIX Daily       & $\text{VIX}_{t}/(100\sqrt{252})$ \\
  & VRP             & Equation~\eqref{eq:vrp} \\
  & Credit Stress   & $\text{HYG}_{t}/\text{IEF}_{t}$ \\
  & Vol Term Slope  & $\text{VIX}_{t}/\text{VIX3M}_{t}$ \\
\midrule
\multirow{2}{*}{Price}
  & FracDiff Close  & $(1-L)^{0.6}\ln P_{t}$; Equation~\eqref{eq:fracdiff} \\
\midrule
\multirow{2}{*}{Regime}
  & Regime 0        & Low-vol GMM dummy (BIC-selected, 2 components) \\
  & Regime 1        & High-vol GMM dummy \\
\bottomrule
\end{tabular}
\end{table}

\subsubsection*{Lookback Window Construction}

At each day $t$, the feature vector presented to the ML models is the
concatenation of 18-dimensional feature observations over the preceding
$L = 20$ trading days:
\begin{equation}
  \mathbf{x}_{t} =
    \bigl[\mathbf{z}_{t-19}^{\top},\,
          \mathbf{z}_{t-18}^{\top},\,\ldots,\,
          \mathbf{z}_{t}^{\top}\bigr]^{\top}
    \in \mathbb{R}^{360},
  \label{eq:lookback}
\end{equation}
where $\mathbf{z}_{s} \in \mathbb{R}^{18}$ is the feature vector at day $s$.
This flattening converts the three-dimensional panel (samples $\times$
features $\times$ lags) into a two-dimensional design matrix that gradient
boosted trees consume directly.  The window of $L = 20$ days corresponds to
approximately one calendar month.  After construction, the training design
matrix has dimensions $1{,}936 \times 360$, and all features are scaled with
a RobustScaler fitted on the training set only, using interquartile range
normalization to limit the influence of the March 2020 volatility spike.

\subsection{Model Specifications}

\subsubsection*{XGBoost}

The base XGBoost model minimizes the squared-error objective on the log-vol
target with $n = 300$ trees, learning rate $\eta = 0.04$, maximum depth 4,
column and row subsampling rates of 0.8, and $L_{2}$ regularization
$\lambda = 2$.

\subsubsection*{LightGBM GBDT}

LightGBM uses standard gradient boosting (GBDT) with $n = 400$ trees,
$\eta = 0.04$, 63 leaves, minimum 20 samples per leaf, $L_{1}$
regularization $\alpha = 0.1$, and $L_{2}$ regularization $\lambda = 2$.
The leaf-wise growth of LightGBM produces trees with up to 63 leaves per
level, compared with XGBoost's maximum of $2^{4} = 16$ leaves under
\texttt{max\_depth} = 4, providing genuine architectural diversity for the
ensemble.  The DART boosting mode was evaluated and discarded: DART produced
out-of-fold MAPE near 77\%, a known instability arising because DART uses all
trees at inference despite the dropout training procedure, leading to
systematic underfitting on regression objectives.

\subsubsection*{Entropy XGBoost (Monotone-Constrained)}

The Entropy XGBoost model augments the XGBoost specification ($n = 350$,
$\eta = 0.03$, $\lambda = 3$) with hard monotone constraints on selected
features.  Let $\mathcal{I}^{+}$ denote the index set of all 20 lookback
lags of the Garman-Klass volatility, Parkinson volatility, and overnight gap
features, and let $\mathcal{I}^{-}$ denote the 20 lags of the log return.
The constraints require:
\begin{align}
  \frac{\partial \hat{f}(\mathbf{x})}{\partial x_{j}} \geq 0
    & \quad \text{for all } j \in \mathcal{I}^{+},
    \label{eq:mono-pos} \\
  \frac{\partial \hat{f}(\mathbf{x})}{\partial x_{j}} \leq 0
    & \quad \text{for all } j \in \mathcal{I}^{-}.
    \label{eq:mono-neg}
\end{align}
These constraints are enforced at the tree-split level: any proposed split
on a feature in $\mathcal{I}^{+}$ that assigns a higher prediction to the
lower-valued branch is rejected, regardless of its improvement to the
training objective.  The constraint in~\eqref{eq:mono-neg} encodes the
leverage effect \citep{Black1976}: higher past returns cannot raise the
volatility forecast.  The constraint in~\eqref{eq:mono-pos} encodes
volatility clustering: higher past realized volatility must produce a higher
forecast.  The remaining 340 feature-lag pairs ($\mathcal{I}^{0}$) are
unconstrained.

\subsubsection*{Quantile Regression}

Three XGBoost models minimize the pinball loss at quantile levels
$\tau \in \{0.10,\, 0.50,\, 0.90\}$:
\begin{equation}
  L_{\tau}(y, q) = (y - q)\bigl[\tau - \mathbf{1}(y < q)\bigr],
  \label{eq:pinball}
\end{equation}
where $y = \text{RV5}_{t+1}$ is the untransformed target (not log-vol).
After fitting, the three forecasts for each observation are sorted by isotonic
regression to enforce $\hat{q}_{0.10} \leq \hat{q}_{0.50} \leq
\hat{q}_{0.90}$, resolving any quantile crossing.  The Q50 conditional median
serves as a point estimator in the main comparison.

\subsubsection*{Stacking Ensembles}

Both stacking configurations use the same out-of-fold (OOF) base model
predictions as meta-features.  The OOF predictions are generated via
five-fold \texttt{TimeSeriesSplit} on the training set with a five-day
embargo: the last five observations of each training fold are excluded before
fitting the base model, preventing label overlap caused by the five shared
return observations between adjacent RV5 windows.  Formally, if fold $k$
has training indices $\mathcal{T}_{k}$ and validation indices
$\mathcal{V}_{k}$, base models are fit on
$\mathcal{T}_{k} \setminus \mathcal{T}_{k}^{\text{end}(5)}$ where
$\mathcal{T}_{k}^{\text{end}(5)}$ denotes the last five indices of
$\mathcal{T}_{k}$.  The resulting OOF matrix has dimensions
$1{,}936 \times 3$.  Scikit-learn's \texttt{StackingRegressor} is not used:
its \texttt{cross\_val\_predict} implementation with \texttt{TimeSeriesSplit}
can assign some observations to multiple validation folds under expanding
windows, producing in-sample predictions that contaminate meta-learner
training.

The Ridge meta-learner is a \texttt{RidgeCV} with $\alpha \in \{0.01, 0.1,
1, 10, 100\}$ fitted on the OOF signal matrix only, without raw feature
pass-through; cross-validation selects $\alpha = 0.1$.  The Huber
meta-learner is a \texttt{HuberRegressor} with $\varepsilon = 1.35$
operating in volatility space (base model log-vol predictions exponentiated
before meta-training) so that the $L_{2}$-to-$L_{1}$ transition at
$1.35\hat{\sigma}$ acts on actual volatility magnitudes rather than
log-volatility residuals.

\subsection{Evaluation Metrics}

Four metrics are reported.  RMSE and MAE measure accuracy in the original
volatility units:
\begin{align}
  \text{RMSE} &= \left(\frac{1}{T}\sum_{t=1}^{T}
    (h_{t} - \sigma_{t})^{2}\right)^{1/2}, \\
  \text{MAE}  &= \frac{1}{T}\sum_{t=1}^{T}|h_{t} - \sigma_{t}|, \\
  \text{MAPE} &= \frac{1}{T}\sum_{t=1}^{T}
    \frac{|h_{t} - \sigma_{t}|}{\sigma_{t}},
\end{align}
where $h_{t}$ is the volatility forecast and $\sigma_{t} = \text{RV5}_{t}$
is the realized volatility proxy.

The primary metric is the QLIKE loss \citep{Patton2011}:
\begin{equation}
  \mathcal{L}_{QLIKE} = \frac{1}{T}\sum_{t=1}^{T}
    \left[\ln H_{t} + \frac{V_{t}}{H_{t}}\right],
  \label{eq:qlike}
\end{equation}
where $H_{t} = h_{t}^{2}$ is the variance forecast and $V_{t} = \sigma_{t}^{2}$
is the realized variance proxy.  The loss is minimized at $H_{t} = V_{t}$,
giving a value of $\ln V_{t} + 1$.  For a typical daily realized variance of
approximately $(0.008)^{2} \approx 6.4 \times 10^{-5}$, the minimum QLIKE is
near $-8.7$, which explains the range of values observed in
Table~\ref{tab:results}.  Forecasts below the perfect-forecast level are
penalized asymmetrically: underpredicting variance ($H_{t} < V_{t}$) raises
$V_{t}/H_{t}$ with diminishing offset from $\ln H_{t}$, while overpredicting
raises $\ln H_{t}$ with diminishing offset from $V_{t}/H_{t}$.

Unlike MSE, QLIKE produces consistent model rankings when $V_{t}$ is a noisy
proxy for latent conditional variance, a property proven in
\citet{Patton2011}.  MSE rankings can be reversed by proxy noise.  This paper
therefore treats QLIKE as the definitive accuracy measure and uses RMSE,
MAE, and MAPE as supporting descriptors.

\subsection{Cross-Validation and Hyperparameter Selection}

Walk-forward cross-validation uses a five-fold \texttt{TimeSeriesSplit} on
the 1,936 training observations, with each fold expanding the training window
from the beginning of the sample.  The five-day embargo is applied at every
fold boundary.  Gradient boosting hyperparameters are fixed at the values
described above; they were selected by RMSE on an internal validation fold
and are not reoptimized on the test set.  The test set is evaluated once at
the end of the study.

% ════════════════════════════════════════════════════════════════════════════
%  5. RESULTS
% ════════════════════════════════════════════════════════════════════════════

\section{Results}
\label{sec:results}

\subsection{Out-of-Sample Performance}

Table~\ref{tab:results} reports RMSE, MAE, QLIKE, and MAPE for all seven
models on the 485-observation test set spanning January 2023 through December
2024.

\begin{table}[htbp]
\centering
\caption{Out-of-Sample Forecast Performance\\
  {\small Test set: January 2023 -- December 2024, $n = 485$ trading days.
   QLIKE is the \citet{Patton2011} quasi-likelihood loss defined in
   Equation~\eqref{eq:qlike}; lower values indicate better accuracy.
   Best performance on each metric is in bold.}}
\label{tab:results}
\begin{tabular}{lcccc}
\toprule
Model & RMSE & MAE & QLIKE & MAPE \\
\midrule
GARCH(1,1)-$t$ & 0.002793 & 0.002167 & $-8.7463$ & 40.89\% \\
XGBoost        & 0.002334 & 0.001770 & $-8.7678$ & 28.26\% \\
LightGBM GBDT  & 0.002441 & 0.001829 & $-8.7429$ & 28.86\% \\
Entropy XGBoost (monotone) & 0.002511 & 0.001885 & $-8.7358$ & 30.53\% \\
Q50 (quantile median) & \textbf{0.002156} & \textbf{0.001607}
                      & $\mathbf{-8.8054}$ & \textbf{25.62\%} \\
Stacking (Ridge meta) & 0.003313 & 0.002547 & $-8.5366$ & 39.61\% \\
Stacking (Huber meta) & 0.003535 & 0.002791 & $-8.4336$ & 48.09\% \\
\bottomrule
\end{tabular}
\end{table}

\paragraph{GARCH(1,1)-$t$.}
The baseline achieves RMSE = 0.002793 and QLIKE $= -8.7463$.  Its MAPE of
40.89\% reflects a persistent overestimation bias in low-volatility periods:
the combination of the positive intercept $\hat{\omega}$ and persistence
$\hat{\alpha} + \hat{\beta}$ close to one causes conditional variance to
decay slowly after spikes, overestimating volatility during the calm periods
that dominate the 2023--2024 test window.

\paragraph{XGBoost.}
XGBoost reduces RMSE by 16.4\% to 0.002334, MAE by 18.3\% to 0.001770, and
MAPE by 12.6 percentage points to 28.26\%.  The QLIKE improvement is
$-8.7678$ versus $-8.7463$, a gain of 0.022.  The modest QLIKE gain relative
to the large RMSE improvement suggests that XGBoost is more accurately
capturing the conditional mean (which drives RMSE) than fine-tuning the tails
of the forecast distribution (which QLIKE measures).

\paragraph{LightGBM GBDT.}
LightGBM achieves RMSE = 0.002441 and MAPE = 28.86\%, closely tracking
XGBoost on scale-sensitive metrics while trailing modestly on QLIKE
($-8.7429$ vs.\ $-8.7678$).  The similar performance of XGBoost and LightGBM
despite their different tree structures (depth-limited vs.\ leaf-wise growth)
indicates that the dominant signal in the feature set is strong enough that
architectural differences in the boosting algorithm matter less than the
quality of the inputs.

\paragraph{Entropy XGBoost.}
The monotone-constrained model achieves RMSE = 0.002511, QLIKE $= -8.7358$,
and MAPE = 30.53\%.  The QLIKE gap relative to unconstrained XGBoost is 0.032,
or approximately 0.4\% relative to the range of QLIKE across models.  This
cost represents the price of the architectural interpretability guarantee:
for any configuration of feature inputs, the Entropy XGBoost model will
never predict that higher recent volatility produces a lower forecast, or that
a positive return shock raises the volatility forecast.

\paragraph{Q50 (Quantile Median).}
The quantile median estimator achieves the best performance on all four
metrics: RMSE = 0.002156, MAE = 0.001607, QLIKE $= -8.8054$, and MAPE =
25.62\%.  Relative to the GARCH baseline, RMSE improves by 22.8\% and QLIKE
improves by 0.059.  Relative to unconstrained XGBoost, RMSE improves by 7.6\%
and QLIKE improves by 0.038.  The Q50 estimator is the single best model in
this study on every reported dimension.

\paragraph{Stacking Ensembles.}
Both stacking configurations underperform every individual base model.  The
Ridge meta-learner achieves RMSE = 0.003313 and QLIKE $= -8.5366$, worse than
GARCH on RMSE by 18.6\% and worse on QLIKE by 0.210.  The Huber meta-learner
achieves RMSE = 0.003535 and QLIKE $= -8.4336$, the worst performance of any
model in the comparison.  These are not marginal underperformances: both
meta-learners produce forecasts that are substantially less accurate than the
individual GARCH baseline, despite being trained on the OOF predictions of
three models that each individually beat GARCH.

\subsection{Quantile Forecast Coverage}

The empirical coverage of the 80\% prediction interval (Q10 to Q90) on the
test set is approximately 80\%, matching the nominal level.  Isotonic
post-processing resolves all quantile crossings and ensures that every
prediction interval has positive width.  The alignment between nominal and
empirical coverage indicates that the quantile models are calibrated rather
than merely conditionally correct.

\subsection{Summary of Relative Performance}

Table~\ref{tab:improvements} expresses each model's performance as a
percentage change in RMSE relative to the GARCH baseline.

\begin{table}[htbp]
\centering
\caption{RMSE Improvement Relative to GARCH(1,1)-$t$ Baseline}
\label{tab:improvements}
\begin{tabular}{lrr}
\toprule
Model & $\Delta$RMSE (\%) & $\Delta$QLIKE (abs.) \\
\midrule
XGBoost                   & $-16.4$ & $-0.022$ (better) \\
LightGBM GBDT             & $-12.6$ & $+0.003$ (worse)  \\
Entropy XGBoost           & $-10.1$ & $+0.011$ (worse)  \\
Q50 (quantile median)     & $-22.8$ & $-0.059$ (better) \\
Stacking (Ridge meta)     & $+18.6$ & $+0.210$ (worse)  \\
Stacking (Huber meta)     & $+26.6$ & $+0.313$ (worse)  \\
\bottomrule
\end{tabular}
\end{table}

% ════════════════════════════════════════════════════════════════════════════
%  6. DISCUSSION
% ════════════════════════════════════════════════════════════════════════════

\section{Discussion}
\label{sec:disc}

\subsection{Why the Quantile Median Estimator Wins}

The Q50 model minimizes the pinball loss at $\tau = 0.5$, which is equivalent
to minimizing the mean absolute error of the forecast.  Its dominance on
RMSE, QLIKE, and MAPE simultaneously suggests that the distribution of
$\text{RV5}_{t+1}$ conditional on $\mathbf{x}_{t}$ is right-skewed with
heavy tails, a feature characteristic of realized volatility series.  Under
such a distribution, the conditional median lies below the conditional mean.
The mean-targeting models (XGBoost, LightGBM, Entropy XGBoost) overpredict
during calm periods and underpredict during spikes because they chase the
mean of a heavy-tailed distribution.  The Q50 model bypasses this by directly
targeting the median, producing forecasts that are better calibrated to the
typical realized volatility observation.

The convergence of best RMSE and best QLIKE in the Q50 model is notable
because RMSE and QLIKE have different sensitivities to the tails of the error
distribution.  The fact that the same model wins on both suggests that median
estimation is not merely robust to outliers at the cost of typical-day
accuracy; it genuinely produces better-calibrated forecasts across the full
distribution.

\subsection{Why Stacking Fails}

The most striking result is the collapse of both stacking configurations.
The most plausible explanation is error correlation among the base models.
All three base models rely on the same dominant feature, lagged realized
volatility, and their forecast errors are highly correlated across
observations.  Stacking gains arise when base model errors are
\emph{complementary}: one model's overestimation offsets another's
underestimation.  When base model errors are aligned, the meta-learner cannot
average away idiosyncratic error; it can only estimate a weighted mean of the
three forecasts, and this mean is worse than the best individual forecast
because the weights are estimated from a noisy 1,936-observation OOF matrix
with only three effective signal dimensions.

A secondary factor is the mismatch between the OOF training environment and
the test environment.  The meta-learners are trained on OOF predictions that
span 2015--2023, a period including the 2015--2016 China volatility shock,
the 2018 VIX spike, and the March 2020 pandemic crash.  The test set covers
2023--2024, a lower-volatility regime.  A meta-learner that learned to combine
base models during high-volatility episodes may not transfer well to a
relatively calm test environment.

This result is consistent with the forecast combination evidence in
\citet{RapachZhou2013}, who find that combination improves on individual
models most reliably when constituent forecasts have complementary information
sets.  In a setting where all models share the same primary predictor, the
combination benefit is limited.

\subsection{Monotone Constraints as an Interpretability Guarantee}

The Entropy XGBoost model represents a specific point on the
interpretability-performance tradeoff curve.  The QLIKE cost of the
constraints is 0.032 relative to unconstrained XGBoost, a difference that is
directionally consistent but may not be statistically significant without a
formal Diebold-Mariano test.  In exchange, the model provides an
architectural guarantee that practitioners can verify without examining
specific trees: no combination of input values can produce a volatility
forecast that contradicts the leverage effect or volatility clustering.

This property is relevant to production risk systems, where a model prediction
that contradicts well-established financial intuition requires either a manual
override or a detailed explanation to risk committees and regulators.  The
constrained model eliminates this failure mode by construction.

\subsection{Limitations}

Several limitations of this study should be acknowledged.

\paragraph{No Diebold-Mariano Test Results.}
The Diebold-Mariano test \citep{DieboldMariano1995} is implemented in the
codebase with Newey-West long-run variance and QLIKE loss but is not applied
to the pairwise model comparisons in Table~\ref{tab:results}.  The QLIKE
differences between, for example, XGBoost and GARCH (0.022) and between Q50
and XGBoost (0.038) are small in absolute terms relative to typical
cross-model variation.  Some of these pairwise differences may not be
statistically significant at conventional levels.  Readers should treat the
rankings in Table~\ref{tab:results} as directional evidence pending formal
significance testing.

\paragraph{HAR-RV Is Not Included.}
The HAR-RV model \citep{Corsi2009} is fully implemented in the codebase and
unit-tested but is not included in the out-of-sample comparison.
\citet{HansenLunde2005} document that HAR-RV consistently outperforms
GARCH(1,1) when realized measures are available as inputs; its exclusion means
this study's ``econometric baseline'' may be weaker than it would be in a
complete comparison.  Including HAR-RV would also clarify whether the ML
models' improvements over GARCH are really improvements over the class of
linear econometric models, or specifically over the GARCH(1,1)-$t$
specification.

\paragraph{No Economic Value Assessment.}
The analysis is purely statistical.  Translating QLIKE improvements into
portfolio performance via a volatility-targeting strategy, where position
sizes are scaled as $k/\hat{h}_{t+1}$, would clarify whether the observed
differences are economically meaningful.  A model that improves RMSE by 22.8\%
but produces economically equivalent Sharpe ratios would lead to different
practical conclusions.

\paragraph{Single Asset and Sample Period.}
The analysis covers one asset (SPY) over one decade.  Results may not
generalize to individual stocks with lower liquidity, different volatility
dynamics, or different information environments.  The test period
(2023--2024) is a relatively low-volatility regime that may favor models
which avoid overestimation, potentially understating the performance advantage
of the GARCH baseline relative to more turbulent periods.

% ════════════════════════════════════════════════════════════════════════════
%  7. CONCLUSION
% ════════════════════════════════════════════════════════════════════════════

\section{Conclusion}
\label{sec:conc}

This paper compared seven volatility forecasting models on a ten-year daily
sample of SPY data using a strict chronological train-test split and QLIKE as
the primary evaluation metric.  Four main findings emerge.

First, GARCH(1,1)-$t$ is a formidable baseline.  Its RMSE of 0.002793 and
QLIKE of $-8.7463$ on the 485-day test set are not easily surpassed, and the
16.4\% RMSE improvement achieved by XGBoost represents meaningful but not
transformative progress.

Second, the quantile median estimator dominates the comparison on every
reported metric, with RMSE = 0.002156 and QLIKE $= -8.8054$.  The result
suggests that the conditional distribution of five-day realized volatility is
sufficiently right-skewed that targeting the conditional median, rather than
the conditional mean, produces better-calibrated forecasts under both symmetric
and asymmetric loss functions.

Third, stacking ensembles fail.  Both meta-learners produce forecasts that are
substantially worse than any of their three component base models.  The failure
is attributed to high correlation among base model errors, which limits the
diversification benefit that makes stacking effective.

Fourth, monotone domain constraints impose a modest QLIKE cost (0.032 relative
to unconstrained XGBoost) in exchange for an architectural guarantee that the
forecast always satisfies the leverage effect and volatility clustering
constraints.  This tradeoff may be favorable in production environments where
model predictions must be defensible to risk committees.

Three extensions would materially strengthen this analysis.  First, applying
the implemented Diebold-Mariano test to all model pairs would convert the
directional rankings in Table~\ref{tab:results} into statements about
statistical significance.  Second, including HAR-RV as an econometric
benchmark would complete the comparison between ML and the class of linear
econometric volatility models.  Third, a volatility-targeting portfolio
backtest would translate the statistical QLIKE improvements into an economic
value assessment.

\newpage

% ════════════════════════════════════════════════════════════════════════════
%  REFERENCES
% ════════════════════════════════════════════════════════════════════════════

\bibliographystyle{apalike}
\bibliography{references}

\end{document}
